% Intro, part after the question (draft, Oct 26, 2026). Not in the mock or the main draft.
% Built from the shipped NERA intro (latex/JMP_DS_nera/sections/01_introduction_v3.tex),
% with "mortgage" removed and section references matched to the mock.
% Every number below was checked against the current mock sections on Oct 26:
%   06_policy.tex (credit 1.3/-1.1, supply 4.35/1.4/3.0/2.0, parents 5.1/16.6/4.9, price 33, households 36,
%   property tax -7.0/+14.5), Table tab:policy-tempo; 05_household_choices.tex (2.9% space cost, risk 12%);
%   02d_motivating_evidence.tex (1.44 rooms, 31%, 27 pp, moves +33%).
% Numbers marked [check] appear in NERA but not in the current mock's policy section.

This paper asks which housing policies raise births, by how much, and over
what horizon. I answer with a quantitative life-cycle model in which
households choose whether to have children, whether to rent or own, how much
space to occupy and how much to save under uninsurable earnings risk, and in
which house prices and rents clear the housing market. I find that policy
raises births when it lowers the cost of a child rather than the cost of a
home, and that its effects build up over a generation. Removing the down
payment for all households raises births by 1.3\% on impact but lowers the
number of children that young adults eventually have by 1.1\%: credit mostly
brings births forward. A housing supply expansion with the same effect on
impact raises that number by 2.0\%, and credit available only to parents
raises it by 4.9\%.

The difference comes from what each policy changes for a household deciding
whether to have a child. Landlords are competitive, so the rent equals the
user cost of owning the same space, and a family can rent its rooms at
roughly the per-period cost of owning them. At the 2007 price, the space a
child requires costs about 2.9\% of mean annual earnings every year. Lowering
the down payment makes it easier to buy a home, but it does not lower this
cost, which the family pays whether it rents or owns. Most of the households
that use it are young couples who buy the small homes they would otherwise
rent. A supply expansion instead lowers prices and rents, and with them the
cost of family space for every cohort that follows. Credit tied to children
rewards the birth itself and lets parents borrow against their home while
their children live with them, when the commitment weighs most.

The model builds on evidence on how housing changes at parenthood
(Section~\ref{sec:empirical-evidence}). In the Panel Study of Income
Dynamics, three to four years after their first birth parents live in 1.44
more rooms than two to three years before it, a 31\% increase, and they are
27 percentage points more likely to own their home. Moves for more space
become a third more likely after the birth, while moves for a better
neighborhood do not. The demand for space jumps when a child arrives, and
owning is the main route to a large home.

In the model (Section~\ref{sec:model}), households live through a life cycle,
save in a bond and face persistent earnings risk that they cannot insure.
During their fertile years they decide whether to try for a child, and a
child raises the household's needs and requires a minimum amount of space for
as long as it lives at home. Households can rent or buy. Rental units are
capped in size, so the largest homes are only for sale, and owners can borrow
up to 80\% of the value of their home while renters cannot borrow, so buying
requires a down payment. Developers build new housing at a cost that rises
with the stock, absentee landlords rent it out, and house prices and rents
clear the market. Children who leave home form households of their own, so
fertility today shapes the demand for housing a generation later.

I calibrate the model to the United States in 2007, before the decline, to
match moments on fertility, housing and wealth, including the rooms parents
add at their first birth (Section~\ref{sec:quant}). I then treat the fall in
births since 2007 as a decline in the preference for children. The model is
not meant to explain the decline; it asks how housing responds to it and what
housing policy can do from now on. Every policy I study starts from the same
2023 economy, with the same households, wealth, homes and children, and is
compared with the path the economy would follow without it.

The results show why the impact of a policy can mislead
(Section~\ref{sec:policy}). The births that credit adds in 2023 fade within a
few years, and births fall below the path without policy in the 2030s,
because many of the children born early would otherwise have been born later.
A supply expansion with the same effect on impact does the opposite: births
keep rising and are 3.0\% above the path without policy by 2063, as each
cohort that enters adulthood spends its fertile years with cheaper space and
the additional children form households of their own. In the terms of
demographers, credit changes the tempo of births and supply changes their
quantum. Credit for parents raises births by 5.1\% in 2023 and by 16.6\% by
2063. A rebated property tax, by contrast, lowers the house price by 7.0\%
but raises rents by 14.5\%, and births barely move: a lower price is not
cheaper space.

Finally, in the long run births per household return to replacement, so the
lasting gains of any policy show up as more households rather than more
children per household. Credit for parents raises fertility at every price,
so the price rises until fertility is back at replacement, 33\% higher in the
new steady state, with 36\% more households. Housing policy aimed at births
should therefore be judged by its long-run effects, not only by its effect on
impact.

% [literature paragraphs go here]

The rest of the paper is organized as follows.
Section~\ref{sec:empirical-evidence} presents the evidence on housing around
the first birth. Section~\ref{sec:model} describes the model, and
Section~\ref{sec:quant} its quantification and the transition from 2007 to
2023. Section~\ref{sec:results} describes household choices in the calibrated
model, Section~\ref{sec:policy} studies housing policies, and
Section~\ref{sec:conclusion} concludes.
